\chapter{2000年大纲卷（理）}
\section{单选题}

\begin{problem}
设集合 \(A\) 和 \(B\) 都是自然数集合 \(\N\)，映射 \(f:A\to B\) 把集合 \(A\) 中的元素 \(n\) 映射成集合 \(B\) 中的元素 \(2^n+n\)，则在映射 \(f\) 下，象 \(20\) 的原象是（\quad）

\choices
    {\(2\)}
    {\(3\)}
    {\(4\)}
    {\(5\)}
\end{problem}

\begin{answer}
C
\end{answer}

\begin{solution}
由 \(f(n)=2^n+n\) 以及 \(f(n+1)-f(n)=2^n+1>0\)，可知 \(f\) 在自然数集合上严格递增. 又 \(f(4)=2^4+4=20\)，所以象 \(20\) 的原象是 \(4\)，故选 C.
\end{solution}

\begin{problem}
在复平面内，把复数 \(3-\sqrt{3}\i\) 对应的向量按顺时针方向旋转 \(\frac{\pi}{3}\)，所得向量对应的复数是（\quad）

\choices
    {\(2\sqrt{3}\)}
    {\(-2\sqrt{3}\i\)}
    {\(\sqrt{3}-3\i\)}
    {\(3+\sqrt{3}\i\)}
\end{problem}

\begin{answer}
B
\end{answer}

\begin{solution}
顺时针旋转 \(\frac{\pi}{3}\) 相当于乘以
\(\cos\frac{\pi}{3}-\i\sin\frac{\pi}{3}=\frac12-\frac{\sqrt{3}}2\i\). 因此所得复数为
\[
(3-\sqrt{3}\i)\left(\frac12-\frac{\sqrt{3}}2\i\right)=-2\sqrt{3}\i.
\]
故选 B.
\end{solution}

\begin{problem}
一个长方体共一顶点的三个面的面积分别是 \(\sqrt{2}\)，\(\sqrt{3}\)，\(\sqrt{6}\)，这个长方体对角线的长是（\quad）

\choices
    {\(2\sqrt{3}\)}
    {\(3\sqrt{2}\)}
    {\(6\)}
    {\(\sqrt{6}\)}
\end{problem}

\begin{answer}
D
\end{answer}

\begin{solution}
设长方体的三条棱长分别为 \(x\)，\(y\)，\(z\)，则
\(xy=\sqrt{2}\)，\(yz=\sqrt{3}\)，\(zx=\sqrt{6}\). 由此得
\(x^2=\frac{(xy)(xz)}{yz}=2\)，\(y^2=\frac{(xy)(yz)}{xz}=1\)，\(z^2=\frac{(xz)(yz)}{xy}=3\). 所以长方体对角线长为
\[
\sqrt{x^2+y^2+z^2}=\sqrt{2+1+3}=\sqrt{6}.
\]
故选 D.
\end{solution}

\begin{problem}
已知 \(\sin\alpha>\sin\beta\)，那么下列命题成立的是（\quad）

\choices
    {若 \(\alpha\)，\(\beta\) 是第一象限角，则 \(\cos\alpha>\cos\beta\)}
    {若 \(\alpha\)，\(\beta\) 是第二象限角，则 \(\tan\alpha>\tan\beta\)}
    {若 \(\alpha\)，\(\beta\) 是第三象限角，则 \(\cos\alpha>\cos\beta\)}
    {若 \(\alpha\)，\(\beta\) 是第四象限角，则 \(\tan\alpha>\tan\beta\)}
\end{problem}

\begin{answer}
D
\end{answer}

\begin{solution}
在第一象限内，正弦函数递增，所以 \(\sin\alpha>\sin\beta\) 推出 \(\alpha>\beta\)，而余弦函数递减，故选项 A 错误. 在第二象限内正弦函数递减，因而 \(\alpha<\beta\)，正切函数递增，选项 B 错误. 在第三象限内正弦函数递减且余弦函数递增，仍有 \(\cos\alpha<\cos\beta\)，选项 C 错误. 在第四象限内正弦函数递增，故 \(\alpha>\beta\)，又正切函数递增，所以 \(\tan\alpha>\tan\beta\). 故选 D.
\end{solution}

\begin{problem}
函数 \(y=-x\cos x\) 的部分图象是（\quad）

\choices
    {\choicebitmap[width=0.15\paperwidth]{img/2000/syllabus_science/q5_optA.png}}
    {\choicebitmap[width=0.15\paperwidth]{img/2000/syllabus_science/q5_optB.png}}
    {\choicebitmap[width=0.15\paperwidth]{img/2000/syllabus_science/q5_optC.png}}
    {\choicebitmap[width=0.15\paperwidth]{img/2000/syllabus_science/q5_optD.png}}
\end{problem}

\begin{answer}
D
\end{answer}

\begin{solution}
函数满足 \(y(0)=0\). 当 \(x\) 在原点右侧且足够接近原点时，\(\cos x>0\)，所以 \(y=-x\cos x<0\)；当 \(x\) 在原点左侧且足够接近原点时，\(y>0\). 此外，\(x=\frac{\pi}{2}\) 是另一个零点，且在 \(0<x<\frac{\pi}{2}\) 时函数值为负，在 \(x>\frac{\pi}{2}\) 且接近该点时函数值为正. 因此图象过原点时由左上方向右下方向穿过，并在右侧先位于 \(x\) 轴下方，符合选项 D. 故选 D.
\end{solution}

\begin{problem}
《中华人民共和国个人所得税法》规定，公民全月工资、薪金所得不超过 \(800\text{ 元}\) 的部分不必纳税，超过 \(800\text{ 元}\) 的部分为全月应纳税所得额，此项税款按下表分段累进计算：

\begin{center}
\begin{tabular}{|c|c|}
\hline
全月应纳税所得额 & 税率 \\
\hline
不超过 \(500\text{ 元}\) 的部分 & \(5\%\) \\
\hline
超过 \(500\text{ 元}\) 至 \(2000\text{ 元}\) 的部分 & \(10\%\) \\
\hline
超过 \(2000\text{ 元}\) 至 \(5000\text{ 元}\) 的部分 & \(15\%\) \\
\hline
\(\cdots\) & \(\cdots\) \\
\hline
\end{tabular}
\end{center}

某人一月份交纳此项税款 \(26.78\text{ 元}\)，则他的当月工资、薪金所得介于（\quad）

\choices
    {\(800\text{ 元}-900\text{ 元}\)}
    {\(900\text{ 元}-1200\text{ 元}\)}
    {\(1200\text{ 元}-1500\text{ 元}\)}
    {\(1500\text{ 元}-2800\text{ 元}\)}
\end{problem}

\begin{answer}
C
\end{answer}

\begin{solution}
设超过免税额的部分为 \(t\) 元，则工资、薪金所得为 \(t+800\) 元. 当 \(0\leqslant t\leqslant500\) 时，税款至多为 \(0.05\times500=25\) 元，不可能是 \(26.78\) 元. 因此 \(500<t\leqslant2000\)，税款为
\[
25+0.1(t-500)=26.78.
\]
解得 \(t=517.8\)，所以当月工资、薪金所得为 \(1317.8\) 元，介于 \(1200\) 元和 \(1500\) 元之间，故选 C.
\end{solution}

\begin{problem}
若 \(a>b>1\)，\(P=\sqrt{\lg a\cdot\lg b}\)，\(Q=\frac12\left(\lg a+\lg b\right)\)，\(R=\lg\left(\frac{a+b}{2}\right)\)，则（\quad）

\choices
    {\(R<P<Q\)}
    {\(P<Q<R\)}
    {\(Q<P<R\)}
    {\(P<R<Q\)}
\end{problem}

\begin{answer}
B
\end{answer}

\begin{solution}
令 \(u=\lg a\)，\(v=\lg b\). 因为 \(a>b>1\)，所以 \(u>v>0\). 由算术平均数不小于几何平均数，且 \(u\ne v\)，得
\[
P=\sqrt{uv}<\frac{u+v}{2}=Q.
\]
又由对数函数的凹性，或由算术平均数与几何平均数的关系，有
\[
R=\lg\left(\frac{a+b}{2}\right)>\frac{\lg a+\lg b}{2}=Q.
\]
因此 \(P<Q<R\)，故选 B.
\end{solution}

\begin{problem}
以极坐标系中的点 \((1,1)\) 为圆心，\(1\) 为半径的圆的方程是（\quad）

\choices
    {\(\rho=2\cos\left(\theta-\frac{\pi}{4}\right)\)}
    {\(\rho=2\sin\left(\theta-\frac{\pi}{4}\right)\)}
    {\(\rho=2\cos\left(\theta-1\right)\)}
    {\(\rho=2\sin\left(\theta-1\right)\)}
\end{problem}

\begin{answer}
C
\end{answer}

\begin{solution}
极坐标点 \((1,1)\) 表示极径为 \(1\)、极角为 \(1\) 的点，圆的半径也为 \(1\). 以极坐标表示的圆心为 \((\rho_0,\theta_0)\)、半径为 \(\rho_0\) 的圆经过极点，其方程为 \(\rho=2\rho_0\cos(\theta-\theta_0)\). 代入 \(\rho_0=1\)，\(\theta_0=1\)，得
\[
\rho=2\cos(\theta-1).
\]
故选 C.
\end{solution}

\begin{problem}
一个圆柱的侧面展开图是一个正方形，这个圆柱的全面积与侧面积的比是（\quad）

\choices
    {\(\frac{1+2\pi}{2\pi}\)}
    {\(\frac{1+4\pi}{4\pi}\)}
    {\(\frac{1+2\pi}{\pi}\)}
    {\(\frac{1+4\pi}{2\pi}\)}
\end{problem}

\begin{answer}
A
\end{answer}

\begin{solution}
设圆柱的底面半径为 \(r\)，高为 \(h\). 侧面展开图是正方形，所以 \(h=2\pi r\). 因此侧面积为 \(2\pi rh=4\pi^2r^2\)，两个底面的面积为 \(2\pi r^2\). 全面积与侧面积的比为
\[
\frac{4\pi^2r^2+2\pi r^2}{4\pi^2r^2}
=\frac{1+2\pi}{2\pi}.
\]
故选 A.
\end{solution}

\begin{problem}
过原点的直线与圆 \(x^2+y^2+4x+3=0\) 相切，若切点在第三象限，则该直线的方程是（\quad）

\choices
    {\(y=\sqrt{3}x\)}
    {\(y=-\sqrt{3}x\)}
    {\(y=\frac{\sqrt{3}}{3}x\)}
    {\(y=-\frac{\sqrt{3}}{3}x\)}
\end{problem}

\begin{answer}
C
\end{answer}

\begin{solution}
将圆方程配方得 \((x+2)^2+y^2=1\)，所以圆心为 \((-2,0)\)，半径为 \(1\). 设过原点的切线为 \(y=kx\)，则圆心到直线的距离等于半径，故
\[
\frac{2|k|}{\sqrt{k^2+1}}=1,
\]
从而 \(k^2=\frac13\)，即 \(k=\pm\frac{\sqrt{3}}3\). 当 \(k=\frac{\sqrt{3}}3\) 时，切点的横、纵坐标均为负，位于第三象限；当 \(k=-\frac{\sqrt{3}}3\) 时切点位于第二象限. 因此所求直线为 \(y=\frac{\sqrt{3}}3x\)，故选 C.
\end{solution}

\begin{problem}
过抛物线 \(y=ax^2\)（\(a>0\)）的焦点 \(F\) 作一直线交抛物线于 \(P\)，\(Q\) 两点，若线段 \(PF\) 与 \(FQ\) 的长分别是 \(p\)，\(q\)，则 \(\frac1p+\frac1q\) 等于（\quad）

\choices
    {\(2a\)}
    {\(\frac1{2a}\)}
    {\(4a\)}
    {\(\frac4a\)}
\end{problem}

\begin{answer}
C
\end{answer}

\begin{solution}
抛物线 \(y=ax^2\) 的焦点为 \(F\left(0,\frac1{4a}\right)\). 设过焦点的直线方程为
\[
y-\frac1{4a}=kx.
\]
它与抛物线的两个交点的横坐标 \(x_1\)，\(x_2\) 是方程
\[
ax^2-kx-\frac1{4a}=0
\]
的两个根. 因此 \(x_1x_2=-\frac1{4a^2}<0\)，两根异号，且
\[
|x_1-x_2|=\frac{\sqrt{k^2+1}}{a}.
\]
由于直线上横坐标相差 \(|x|\) 时距离相差 \(|x|\sqrt{1+k^2}\)，且 \(x_1\)，\(x_2\) 异号，所以 \(|x_1|+|x_2|=|x_1-x_2|\). 因而
\[
\frac1p+\frac1q
=\frac1{\sqrt{1+k^2}}\left(\frac1{|x_1|}+\frac1{|x_2|}\right)
=\frac{|x_1-x_2|}{|x_1x_2|\sqrt{1+k^2}}
=4a.
\]
故选 C.
\end{solution}

\begin{problem}
如图，\(OA\) 是圆锥底面中心 \(O\) 到母线的垂线，\(OA\) 绕轴旋转一周所得曲面将圆锥分成体积相等的两部分，则母线与轴的夹角为（\quad）

\FigureLayoutDeclare{img/2000/syllabus_science/q12_fig1.png}{4.8cm}{0bp 0bp 35bp 0bp}
\bitmapfigure[max width=0.55\linewidth]{img/2000/syllabus_science/q12_fig1.png}

\choices
    {\(\arccos\frac1{\sqrt[3]{2}}\)}
    {\(\arccos\frac12\)}
    {\(\arccos\frac1{\sqrt{2}}\)}
    {\(\arccos\frac1{\sqrt[4]{2}}\)}
\end{problem}

\begin{answer}
D
\end{answer}

\begin{solution}
设圆锥的高为 \(h\)，底面半径为 \(r\)，顶点为 \(V\)，并在含轴的截面内设底面圆周上的母线端点为 \(B\). 取轴为纵轴，令 \(V=(0,h)\)，\(O=(0,0)\)，\(B=(r,0)\). 设 \(A\) 为 \(O\) 到母线 \(VB\) 的垂足，记 \(A\) 到轴的距离为 \(x_A\). 直线 \(VB\) 的方程为 \(hx+ry=hr\)，由垂足坐标公式得
\[
x_A=\frac{rh^2}{r^2+h^2}.
\]
若 \(\theta\) 为母线与轴的夹角，则 \(\cos\theta=\frac{h}{\sqrt{r^2+h^2}}\)，所以 \(x_A=r\cos^2\theta\).

旋转线段 \(OA\) 所得的曲面把原圆锥分成两个部分，其中含有轴截面三角形 \(VOA\) 的一部分由两个底面半径为 \(x_A\)、高分别为 \(y_A\) 和 \(h-y_A\) 的圆锥组成，其中 \(y_A\) 为点 \(A\) 的纵坐标. 因此该部分体积为
\[
\frac13\pi x_A^2y_A+\frac13\pi x_A^2(h-y_A)=\frac13\pi h x_A^2.
\]
原圆锥体积为 \(\frac13\pi h r^2\). 两部分体积相等，因此
\[
\frac{\frac13\pi h x_A^2}{\frac13\pi h r^2}=\frac12,
\]
从而 \(\frac{x_A}{r}=\frac1{\sqrt2}\). 于是 \(\cos^2\theta=\frac1{\sqrt2}\)，且 \(\theta\) 为锐角，故
\[
\cos\theta=\frac1{\sqrt[4]{2}}.
\]
故选 D.
\end{solution}

\section{填空题}

\begin{problem}
乒乓球队的 \(10\) 名队员中有 \(3\) 名主力队员，派 \(5\) 名参加比赛，\(3\) 名主力队员要安排在第一、三、五位置，其余 \(7\) 名队员选 \(2\) 名安排在第二、四位置，那么不同的出场安排共有\fillinblank{}种（用数字作答）.
\end{problem}

\begin{answer}
252
\end{answer}

\begin{solution}
先把 \(3\) 名主力队员安排在第一、三、五位置，有 \(3!\) 种安排. 再从其余 \(7\) 名队员中选出 \(2\) 名，并把他们安排在第二、四位置，有 \(\mathrm C_7^2\cdot2!\) 种安排. 因此不同的出场安排共有
\[
3!\,\mathrm C_7^2\,2!=6\times21\times2=252
\]
种.
\end{solution}

\begin{problem}
椭圆 \(\frac{x^2}{9}+\frac{y^2}{4}=1\) 的焦点 \(F_1\)，\(F_2\)，点 \(P\) 为其上的动点，当 \(\angle F_1PF_2\) 为钝角时，点 \(P\) 横坐标的取值范围是\fillinblank{}.
\end{problem}

\begin{answer}
\(-\frac{3}{\sqrt{5}}<x<\frac{3}{\sqrt{5}}\)
\end{answer}

\begin{solution}
椭圆的长半轴为 \(3\)，短半轴为 \(2\)，所以 \(c=\sqrt{3^2-2^2}=\sqrt5\)，可设
\(F_1=(-\sqrt5,0)\)，\(F_2=(\sqrt5,0)\)，\(P=(x,y)\). 角 \(\angle F_1PF_2\) 为钝角，当且仅当
\[
\overrightarrow{PF_1}\cdot\overrightarrow{PF_2}<0.
\]
而
\[
\overrightarrow{PF_1}\cdot\overrightarrow{PF_2}
=(-\sqrt5-x)(\sqrt5-x)+y^2=x^2+y^2-5.
\]
因此钝角条件为 \(x^2+y^2<5\). 由椭圆方程得 \(y^2=4-\frac49x^2\)，代入可得
\[
4+\frac59x^2<5,
\]
即 \(x^2<\frac95\). 所以点 \(P\) 的横坐标范围为
\[
-\frac3{\sqrt5}<x<\frac3{\sqrt5}.
\]
\end{solution}

\begin{problem}
设 \(\{a_n\}\) 是首项为 \(1\) 的正项数列，且 \((n+1)a_{n+1}^2-na_n^2+a_{n+1}a_n=0\)（\(n=1,2,3,\dots\)），则它的通项公式是 \(a_n=\fillinblank{}\).
\end{problem}

\begin{answer}
\(\frac{1}{n}\)
\end{answer}

\begin{solution}
因为数列各项为正，令 \(r_n=\frac{a_{n+1}}{a_n}>0\). 将递推关系除以 \(a_n^2\)，得
\[
(n+1)r_n^2+r_n-n=0.
\]
这个关于 \(r_n\) 的方程判别式为 \(1+4n(n+1)=(2n+1)^2\)，所以
\[
r_n=\frac{-1+(2n+1)}{2(n+1)}=\frac{n}{n+1},
\]
另一个根为负数，不符合 \(r_n>0\). 于是
\[
a_{n+1}=\frac{n}{n+1}a_n.
\]
由 \(a_1=1\) 逐项相乘，得到
\[
a_n=a_1\prod_{k=1}^{n-1}\frac{k}{k+1}
=\frac1n.
\]
故通项公式为 \(a_n=\frac1n\).
\end{solution}

\begin{problem}
如图，\(E\)，\(F\) 分别为正方体面 \(ADD_1A_1\)，面 \(BCC_1B_1\) 的中心，则四边形 \(BFD_1E\) 在该正方体的面上的射影可能是\fillinblank{}（要求：把可能的图序号都填上）.

\begingroup
\leftskip=0pt\relax
\everypar{}
\hbadness=10000
\examfiguregroup[float=inline]{img/2000/syllabus_science/q16_fig1.png}{%
    \bitmapinclude[max width=0.82\linewidth]{img/2000/syllabus_science/q16_fig1.png}\par
    \noindent
    \makebox[0.22\linewidth][c]{\bitmapinclude[width=0.20\linewidth]{img/2000/syllabus_science/q16_optA.png}}\hfill
    \makebox[0.22\linewidth][c]{\bitmapinclude[width=0.20\linewidth]{img/2000/syllabus_science/q16_optB.png}}\hfill
    \makebox[0.22\linewidth][c]{\bitmapinclude[width=0.20\linewidth]{img/2000/syllabus_science/q16_optC.png}}\hfill
    \makebox[0.22\linewidth][c]{\bitmapinclude[width=0.20\linewidth]{img/2000/syllabus_science/q16_optD.png}}%
}
\endgroup
\end{problem}

\begin{answer}
\(\circled{2}\circled{3}\)
\end{answer}

\begin{solution}
取正方体棱长为 \(1\)，设
\[
A=(0,0,0),\quad B=(1,0,0),\quad D_1=(0,1,1),\quad
E=\left(0,\frac12,\frac12\right),\quad
F=\left(1,\frac12,\frac12\right).
\]
在两个与面 \(ADD_1A_1\)、\(BCC_1B_1\) 平行的面上作射影时，\(E\) 与 \(F\) 的射影重合，\(B\) 与 \(D_1\) 的射影为该面的两个对角顶点，所以四边形的射影退化为一条对角线，对应图 \(\circled{3}\).

在其余四个面上作射影时，以相应面的两条棱为坐标轴，四个顶点的坐标可写成
\[
\left(1,0\right),\quad\left(1,\frac12\right),\quad
\left(0,1\right),\quad\left(0,\frac12\right).
\]
所得图形为图 \(\circled{2}\) 所示的四边形. 因此可能的图序号为 \(\circled{2}\)、\(\circled{3}\).
\end{solution}

\section{解答题}

\begin{problem}
已知函数 \(y=\sqrt{3}\sin x+\cos x\)（\(x\in\R\)）.

\begin{enumerate}
\item 当函数 \(y\) 取得最大值时，求自变量 \(x\) 的集合.
\item 该函数的图象可由 \(y=\sin x\)（\(x\in\R\)）的图象经过怎样的平移和伸缩变换得到？
\end{enumerate}
\end{problem}

\begin{solution}
\begin{enumerate}
\item 利用两角和公式，
\[
\sqrt{3}\sin x+\cos x
=2\left(\frac{\sqrt{3}}{2}\sin x+\frac{1}{2}\cos x\right)
=2\sin\left(x+\frac{\pi}{6}\right).
\]
由于正弦函数的最大值为 \(1\)，所以函数 \(y\) 的最大值为 \(2\). 当且仅当
\[
x+\frac{\pi}{6}=\frac{\pi}{2}+2k\pi\quad (k\in\Z)
\]
时取到最大值，故自变量 \(x\) 的集合为
\[
\left\{\frac{\pi}{3}+2k\pi\mid k\in\Z\right\}.
\]
\item 先将 \(y=\sin x\) 的图象上各点的纵坐标伸长为原来的 \(2\) 倍，得到 \(y=2\sin x\) 的图象；再将所得图象向左平移 \(\frac{\pi}{6}\)，得到
\[
y=2\sin\left(x+\frac{\pi}{6}\right)=\sqrt{3}\sin x+\cos x.
\]
因此，所求图象可由 \(y=\sin x\) 的图象先作纵向伸长，再向左平移 \(\frac{\pi}{6}\) 得到.
\end{enumerate}
\end{solution}

\begin{problem}
如图，已知平行六面体 \(ABCD-A_1B_1C_1D_1\) 的底面 \(ABCD\) 为菱形，且 \(\angle C_1CB=\angle C_1CD=\angle BCD=60^\circ\).

\begin{enumerate}
\item 证明：\(C_1C\perp BD\)；
\item 假定 \(CD=2\)，\(CC_1=\frac32\)，记面 \(C_1BD\) 为 \(\alpha\)，面 \(CBD\) 为 \(\beta\)，求二面角 \(\alpha-BD-\beta\) 的平面角的余弦值；
\item 当 \(\frac{CD}{CC_1}\) 的值为多少时，能使 \(A_1C\perp\) 平面 \(C_1BD\)？请给出证明.
\end{enumerate}

\bitmapfigure[max width=0.7\linewidth]{img/2000/syllabus_science/q18_fig1.png}
\end{problem}

\begin{solution}
\begin{enumerate}
\item 设 \(\bs{b}=\overrightarrow{CB}\)，\(\bs{d}=\overrightarrow{CD}\)，\(\bs{c}=\overrightarrow{CC_1}\). 由已知角相等，
\[
\bs{c}\cdot\bs{b}=|\bs{c}|\,|\bs{b}|\cos60^\circ=\bs{c}\cdot\bs{d}.
\]
又 \(\overrightarrow{BD}=\bs{d}-\bs{b}\)，所以
\[
\bs{c}\cdot\overrightarrow{BD}=\bs{c}\cdot(\bs{d}-\bs{b})=0.
\]
因此 \(C_1C\perp BD\).

\item 由于 \(CB=CD=2\)，且 \(\angle BCD=60^\circ\)，所以三角形 \(BCD\) 为边长 \(2\) 的等边三角形. 设 \(M\) 为 \(BD\) 的中点，则
\[
CM\perp BD,\qquad CM=\sqrt{3}.
\]
由 \(M\) 为 \(BD\) 的中点，
\[
\overrightarrow{CM}=\frac{\bs{b}+\bs{d}}{2},\qquad
\bs{c}\cdot\overrightarrow{CM}=\frac{\bs{c}\cdot\bs{b}+\bs{c}\cdot\bs{d}}{2}=\frac{3}{2}.
\]
于是
\[
\cos\angle C_1CM=\frac{\bs{c}\cdot\overrightarrow{CM}}{CC_1\cdot CM}
=\frac{3/2}{(3/2)\sqrt{3}}=\frac{\sqrt{3}}{3}.
\]
由于 \(C_1C\perp BD\) 且 \(CM\perp BD\)，直线 \(MC_1\) 也垂直于 \(BD\)，所以二面角 \(\alpha-BD-\beta\) 的平面角为 \(\angle C_1MC\). 在三角形 \(C_1CM\) 中，
\[
C_1M^2=CC_1^2+CM^2-2\cdot CC_1\cdot CM\cos\angle C_1CM
=\frac94+3-3=\frac94,
\]
故 \(C_1M=\frac32=CC_1\). 从而
\[
\cos\angle C_1MC=\frac{C_1M^2+CM^2-CC_1^2}{2\cdot C_1M\cdot CM}
=\frac{3}{3\sqrt{3}}=\frac{\sqrt{3}}{3}.
\]

\item 设 \(CD=CB=s>0\)，\(CC_1=h>0\)，仍取 \(M\) 为 \(BD\) 的中点. 由等边三角形 \(BCD\)，
\[
CM=\frac{\sqrt{3}}{2}s,\qquad CM^2=\frac34s^2,
\]
并且
\[
\bs{c}\cdot\overrightarrow{CM}=\frac{hs}{2}.
\]
由平行六面体的向量关系，
\[
\overrightarrow{CA_1}=\bs{b}+\bs{d}+\bs{c}=2\overrightarrow{CM}+\bs{c}.
\]
由第一问，\(\bs{c}\perp BD\). 又因底面 \(ABCD\) 为菱形，\(|\bs{b}|=|\bs{d}|\)，所以
\[
(\bs{b}+\bs{d})\cdot(\bs{d}-\bs{b})=|\bs{d}|^2-|\bs{b}|^2=0.
\]
因此 \(\bs{b}+\bs{d}\perp BD\)，从而
\[
\overrightarrow{CA_1}=\bs{b}+\bs{d}+\bs{c}\perp BD.
\]
平面 \(C_1BD\) 内过 \(M\) 的两条相交直线可取 \(BD\) 与 \(MC_1\)，其中
\[
\overrightarrow{MC_1}=\bs{c}-\overrightarrow{CM}.
\]
因此，\(A_1C\perp\) 平面 \(C_1BD\) 的充要条件是
\[
\overrightarrow{CA_1}\cdot\overrightarrow{MC_1}=0.
\]
代入上式得
\[
\overrightarrow{CA_1}\cdot\overrightarrow{MC_1}
=(2\overrightarrow{CM}+\bs{c})\cdot(\bs{c}-\overrightarrow{CM})
=h^2+\frac{hs}{2}-\frac32s^2.
\]
所以
\[
2h^2+hs-3s^2=0.
\]
令 \(r=\frac{s}{h}>0\)，则
\[
3r^2-r-2=0,\qquad (3r+2)(r-1)=0.
\]
由于 \(r>0\)，只能有 \(r=1\)，即
\[
\frac{CD}{CC_1}=1.
\]
反之，当该比值为 \(1\) 时上面的内积为 \(0\)，且 \(\overrightarrow{CA_1}\perp BD\)，故确实有 \(A_1C\perp\) 平面 \(C_1BD\).
\end{enumerate}
\end{solution}

\begin{problem}
设函数 \(f(x)=\sqrt{x^2+1}-ax\)，其中 \(a>0\).

\begin{enumerate}
\item 解不等式 \(f(x)\leqslant1\)；
\item 求 \(a\) 的取值范围，使函数 \(f(x)\) 在区间 \([0,+\infty)\) 上是单调函数.
\end{enumerate}
\end{problem}

\begin{solution}
\begin{enumerate}
\item 原不等式等价于
\[
\sqrt{x^2+1}\leqslant 1+ax.
\]
由于左边非负，必须有 \(1+ax\geqslant0\). 在此条件下两边平方，得
\[
x^2+1\leqslant(1+ax)^2,
\]
即
\[
x\bigl((1-a^2)x-2a\bigr)\leqslant0.
\]
\(0<a<1\) 时，第二个不等式的两个根为 \(0\) 和 \(\frac{2a}{1-a^2}>0\)，故其解集为 \(\left[0,\frac{2a}{1-a^2}\right]\)，且其中每个 \(x\) 都满足 \(1+ax>0\). 因此此时原不等式的解集为
\[
\left[0,\frac{2a}{1-a^2}\right].
\]
当 \(a=1\) 时，平方后的不等式为 \(-2x\leqslant0\)，结合 \(1+x\geqslant0\)，得到解集 \([0,+\infty)\). 当 \(a>1\) 时，平方后的不等式解为
\[
\left(-\infty,-\frac{2a}{a^2-1}\right]\cup[0,+\infty).
\]
但 \(x\leqslant-\frac{2a}{a^2-1}\) 时有 \(1+ax<0\)，不符合平方前的必要条件，所以原不等式的解集仍为 \([0,+\infty)\). 综上，
\[
\begin{cases}
\left[0,\dfrac{2a}{1-a^2}\right],&0<a<1,\\
[0,+\infty),&a\geqslant1.
\end{cases}
\]
\item 在 \(x\geqslant0\) 时，函数的导数为
\[
f'(x)=\frac{x}{\sqrt{x^2+1}}-a.
\]
又因为 \(0\leqslant\frac{x}{\sqrt{x^2+1}}<1\). 当 \(a\geqslant1\) 时，\(f'(x)<0\)，所以 \(f(x)\) 在 \([0,+\infty)\) 上单调递减.

当 \(0<a<1\) 时，令
\[
x_0=\frac{a}{\sqrt{1-a^2}}>0.
\]
由 \(\frac{x_0}{\sqrt{x_0^2+1}}=a\)，可知 \(f'(x)<0\)（\(0\leqslant x<x_0\)），而 \(f'(x)>0\)（\(x>x_0\)）. 因此函数先减后增，不是单调函数. 故所求 \(a\) 的取值范围为
\[
a\geqslant1.
\]
\end{enumerate}
\end{solution}

\begin{problem}
\begin{enumerate}
\item 已知数列 \(\{c_n\}\)，其中 \(c_n=2^n+3^n\)，且数列 \(\{c_{n+1}-pc_n\}\) 为等比数列，求常数 \(p\)；
\item 设 \(\{a_n\}\)，\(\{b_n\}\) 是公比不相等的两个等比数列，\(c_n=a_n+b_n\)，证明数列 \(\{c_n\}\) 不是等比数列.
\end{enumerate}
\end{problem}

\begin{solution}
\begin{enumerate}
\item 记
\[
d_n=c_{n+1}-pc_n=(2-p)2^n+(3-p)3^n.
\]
若 \(\{d_n\}\) 为等比数列，则必满足 \(d_{n+1}^2=d_nd_{n+2}\). 令 \(A=2-p\)，\(B=3-p\)，则
\[
d_{n+1}^2-d_nd_{n+2}
=(2A2^n+3B3^n)^2-(A2^n+B3^n)(4A2^n+9B3^n)
=-AB6^n=-(2-p)(3-p)6^n.
\]
由于 \(6^n\ne0\)，必须有 \((2-p)(3-p)=0\)，所以 \(p=2\) 或 \(p=3\). 当 \(p=2\) 时 \(d_n=3^n\)，当 \(p=3\) 时 \(d_n=-2^n\)，二者均为等比数列，故这两个值都成立.
\item 设 \(\{a_n\}\) 和 \(\{b_n\}\) 的公比分别为 \(r\) 和 \(s\)，其中 \(r\ne s\). 则
\[
c_{n+1}=ra_n+sb_n,\qquad c_{n+2}=r^2a_n+s^2b_n.
\]
于是
\[
c_{n+1}^2-c_nc_{n+2}
=(ra_n+sb_n)^2-(a_n+b_n)(r^2a_n+s^2b_n)
=-(r-s)^2a_nb_n.
\]
由于两个等比数列的各项非零，且 \(r\ne s\)，右端不为 \(0\). 但是任一等比数列都满足 \(c_{n+1}^2=c_nc_{n+2}\)，矛盾. 因此 \(\{c_n\}\) 不是等比数列.
\end{enumerate}
\end{solution}

\begin{problem}
某蔬菜基地种植西红柿，由历年市场行情得知，从二月一日起的 \(300\) 天内，西红柿市场售价与上市时间的关系用图一的一条折线表示；西红柿的种植成本与上市时间的关系用图二的抛物线段表示.

\begin{enumerate}
\item 写出图一表示的市场售价与时间的函数关系式 \(p=f(t)\)；写出图二表示的种植成本与时间的函数关系式 \(Q=g(t)\)；
\item 认定市场售价减去种植成本为纯收益，问何时上市的西红柿纯收益最大？
\end{enumerate}

注：市场售价和种植成本的单位：\(\text{ 元}/10^{2}\mathrm{~kg}\)，时间单位：\(\text{天}\).

\bitmapinlinefigure[0.75\linewidth][max width=0.75\linewidth]{img/2000/syllabus_science/q21_fig1.png}
\end{problem}

\begin{solution}
\begin{enumerate}
\item 由图一，市场售价折线经过 \((0,300)\)、\((200,100)\)、\((300,300)\). 前一段斜率为 \(-1\)，后一段斜率为 \(2\)，所以
\[
p=f(t)=\begin{cases}
300-t,&0\leqslant t\leqslant200,\\
2t-300,&200<t\leqslant300.
\end{cases}
\]
由图二，成本曲线的顶点为 \((150,100)\)，可设
\[
Q=g(t)=k(t-150)^2+100.
\]
又曲线经过 \((50,150)\)，代入得
\[
150=100+k(50-150)^2=100+10000k,
\]
所以 \(k=\frac{1}{200}\)，从而
\[
g(t)=\frac{1}{200}(t-150)^2+100\quad(0\leqslant t\leqslant300).
\]
\item 设纯收益为 \(H(t)=f(t)-g(t)\). 当 \(0\leqslant t\leqslant200\) 时，
\[
H(t)=300-t-\left[\frac{1}{200}(t-150)^2+100\right]
=100-\frac{(t-50)^2}{200}.
\]
因此这一段的最大值在 \(t=50\) 时取得，最大值为 \(100\).

当 \(200<t\leqslant300\) 时，
\[
H(t)=2t-300-\left[\frac{1}{200}(t-150)^2+100\right]
=100-\frac{(t-350)^2}{200}.
\]
在区间 \((200,300]\) 内，\(t\) 越大，\((t-350)^2\) 越小，所以这一段的最大值在 \(t=300\) 时取得，等于 \(\frac{175}{2}\). 因为 \(100>\frac{175}{2}\)，故纯收益最大时 \(t=50\)，即从二月一日起第 \(50\) 天上市.
\end{enumerate}
\end{solution}

\begin{problem}
如图，已知梯形 \(ABCD\) 中 \(|AB|=2|CD|\)，点 \(E\) 分有向线段 \(\overrightarrow{AC}\) 所成的比为 \(\lambda\)，双曲线过 \(C\)，\(D\)，\(E\) 三点，且以 \(A\)，\(B\) 为焦点，当 \(\frac23\leqslant\lambda\leqslant\frac34\) 时，求双曲线离心率 \(e\) 的取值范围.

\bitmapfigure[max width=0.55\linewidth]{img/2000/syllabus_science/q22_fig1.png}
\end{problem}

\begin{solution}
设双曲线的中心为 \(O\)，焦点 \(A\)，\(B\) 分别为 \((-c,0)\)，\((c,0)\)，则 \(AB=2c\)，且双曲线方程可写为
\[
\frac{x^2}{a^2}-\frac{y^2}{b^2}=1,\qquad c^2=a^2+b^2,\qquad e=\frac{c}{a}.
\]
由图中 \(CD\parallel AB\)，且 \(C,D\) 在双曲线的两支上，设它们的纵坐标均为 \(y_0\). 双曲线方程关于 \(y\) 轴对称，所以 \(C,D\) 的横坐标互为相反数. 又
\[
CD=\frac{AB}{2}=c,
\]
故可设
\[
C\left(\frac{c}{2},y_0\right),\qquad D\left(-\frac{c}{2},y_0\right).
\]
由点 \(C\) 在双曲线上，
\[
\frac{c^2}{4a^2}-\frac{y_0^2}{b^2}=1.
\]
由 \(\frac{AE}{EC}=\lambda\)，点 \(E\) 的坐标为内分点公式
\[
E=\frac{A+\lambda C}{1+\lambda}
=\left(\frac{c(\lambda-2)}{2(1+\lambda)},\frac{\lambda y_0}{1+\lambda}\right).
\]
因为 \(E\) 也在双曲线上，
\[
\frac{c^2(\lambda-2)^2}{4a^2(1+\lambda)^2}
-\frac{\lambda^2y_0^2}{b^2(1+\lambda)^2}=1.
\]
令 \(e^2=\frac{c^2}{a^2}\)，由点 \(C\) 的关系得到
\[
\frac{y_0^2}{b^2}=\frac{e^2}{4}-1.
\]
将其代入点 \(E\) 的关系，并乘以 \((1+\lambda)^2\)，得
\[
\frac{e^2(\lambda-2)^2}{4}-\lambda^2\left(\frac{e^2}{4}-1\right)=(1+\lambda)^2.
\]
化简可得
\[
e^2(1-\lambda)+\lambda^2=1+2\lambda+\lambda^2,\qquad
e^2=\frac{1+2\lambda}{1-\lambda}.
\]
由于 \(\frac23\leqslant\lambda\leqslant\frac34<1\)，且
\[
\frac{\mathrm d}{\mathrm d\lambda}\left(\frac{1+2\lambda}{1-\lambda}\right)=\frac{3}{(1-\lambda)^2}>0,
\]
所以 \(e^2\) 随 \(\lambda\) 增大而增大. 因此
\[
e^2\in\left[\frac{1+2\cdot\frac23}{1-\frac23},\frac{1+2\cdot\frac34}{1-\frac34}\right]=[7,10].
\]
由 \(e>0\)，得到
\[
\sqrt{7}\leqslant e\leqslant\sqrt{10}.
\]
\end{solution}
